% Historical independent verification before the MPP2025 source was located.
% This is a recheck of a known localized statement, not a new theorem claim.
\begin{proposition}[K17-ST2-LOCAL]\label{prop:steinberg-two-local}
For every $n\ge3$, put $R_n=\mathbb Z[1/(n+1)!]$ and let
$\mathrm{St}_n=\widetilde H_{n-2}(\mathcal T_n(\mathbb Q),\mathbb Z)$
be the integral Steinberg module. Then
\[
 H_2(\mathrm{SL}_n(\mathbb Z);\mathrm{St}_n\otimes R_n)=0,
 \qquad
 H_2(\mathrm{GL}_n(\mathbb Z);\mathrm{St}_n\otimes R_n)=0.
\]
\end{proposition}

\begin{proof}
Put $C_i=C_{n-1+i}(\mathrm{BAA}_n,\mathrm{BAA}'_n)$ for
$i=0,1,2$. Br\"uck--Miller--Patzt--Sroka--Wilson construct an integral
exact sequence
\[
 C_2\longrightarrow C_1\longrightarrow C_0
 \longrightarrow\mathrm{St}_n\longrightarrow0.
\]
in their Proposition~8.5, author preprint p.~57
\cite{BrueckEtAl2024}. Localization is flat, so this sequence stays exact
over $R_n$. Their Lemma~8.7, p.~58, says that all three chain modules
are projective over $R_n[\mathrm{SL}_n(\mathbb Z)]$ once the orders of
all finite-order elements of the group are invertible in $R_n$.
If a prime $q$ divides the order of such an element, some power has
order $q$. Its characteristic polynomial over $\mathbb Q$ has the
cyclotomic polynomial $\Phi_q$ as a factor, since an order-$q$ matrix
in characteristic zero is diagonalizable and is not the identity.
Thus $q-1\le n$. Every relevant order is therefore a unit in $R_n$.

The proof of their Lemma~8.9, pp.~58--59, supplies, for every
oriented top simplex, an element of $\mathrm{SL}_n(\mathbb Z)$
preserving that simplex and reversing its orientation. Its two
explicit cases are the double-double and double-triple simplices.
Consequently every generator of the top coinvariant module equals
its own negative. Since $2$ is a unit in $R_n$,
\[
 C_{n+1}(\mathrm{BAA}_n,\mathrm{BAA}'_n;R_n)
     \otimes_{R_n[\mathrm{SL}_n(\mathbb Z)]}R_n=0.
\]
Extend the displayed partial projective resolution to a projective
resolution. Its degree-two coinvariant chain module is zero, and
therefore its degree-two homology is zero. This proves the first
assertion. For the extension
$1\to\mathrm{SL}_n(\mathbb Z)\to\mathrm{GL}_n(\mathbb Z)
\to C_2\to1$, the homology spectral sequence collapses to
coinvariants because $2$ is invertible. In degree two it identifies
the required $\mathrm{GL}_n$ group with the coinvariants of the zero
$\mathrm{SL}_n$ group. This proves the second assertion.
\end{proof}

The new statement is the localized integral consequence of the cited
resolution and its explicit orientation argument. We do not claim a
new proof of the resolution or its computer-assisted connectivity
input. In particular, at $n=11$ it gives the vanishing of the
rank-eleven term after inverting $12!$.

